\chapter{Theoretical Background: Autonomous Web Spider}
\label{chap:theoretical_background_web_spider}

\section{Autonomous Web Data Extraction}
Autonomous web data extraction is a response to a practical interoperability gap: institutions frequently publish operationally relevant information (timetables, announcements, room updates, event schedules) on web pages that are not supported by stable, documented APIs. In such environments, the web itself becomes the integration layer, and a web spider provides a controlled method for discovering and synchronizing content into a structured platform \cite{Sommerville2010}. Unlike ad-hoc scripts that target a single page, an autonomous spider is designed as a repeatable pipeline: it discovers candidate URLs, fetches documents under clear constraints, extracts structured fields, validates outputs, and continuously monitors change over time.

From a theoretical perspective, the spider’s behavior can be modeled as a traversal problem over a directed graph \(G=(V,E)\), where pages are nodes and hyperlinks are edges. The spider begins from one or more seed URLs that anchor the crawl to relevant domains. It then expands the crawl frontier using an explicit traversal policy, subject to constraints such as maximum depth, domain allowlists, fetch-rate limits, and deduplication rules \cite{Narayan2012}. The key design challenge is balancing coverage and cost: the spider must discover sufficiently many relevant pages while avoiding unbounded exploration and unnecessary requests.

\subsection{Traversal Strategies and Breadth-First Search (BFS)}
Traversal strategy determines the order in which pages are visited and therefore shapes discovery quality, latency, and resource consumption. Breadth-First Search (BFS) explores the graph level by level. Concretely, BFS maintains a queue (frontier) and processes all pages at distance \(d\) from the seeds before proceeding to distance \(d+1\). This is valuable in institutional settings because relevant content is often located close to well-known root pages (faculty news pages, course listings, schedules, and category landing pages). BFS therefore tends to discover near-root pages early, offering predictable coverage under a depth limit \cite{Sommerville2010}.

Figure~\ref{fig:bfs_spider} illustrates the BFS intuition: seeds are inserted into the frontier first (distance 0), then discovered links are appended and processed in increasing distance order. A visited set prevents duplicate processing and reduces the risk of infinite loops (e.g., calendars with next/previous navigation). In practice, BFS is frequently combined with additional prioritization heuristics, such as preferring URLs that match known patterns (e.g., \texttt{/news/}, \texttt{/schedule/}) or deprioritizing assets and non-content pages.

In formal terms, the web can be modeled as a directed graph where pages are nodes and hyperlinks are edges. Breadth-First Search (BFS) is a traversal strategy that explores this graph level by level, beginning from one or more seed URLs. BFS maintains a queue of frontier nodes, visits pages in increasing distance from the seed set, and records discovered links for subsequent expansion. Compared with depth-first traversal, BFS tends to discover broad thematic coverage earlier and is suitable when the extractor aims to find recent announcements distributed across many sections of a site. It also provides a natural basis for controlling crawl depth, prioritizing domain boundaries, and preventing unbounded exploration through visited sets and URL normalization rules \cite{Sommerville2010}.

The traversal logic used by a spider is illustrated in Figure~\ref{fig:bfs_spider}. The diagram should be read from left to right. The first box contains the seed URLs (initial sources). The second box represents the active BFS queue (distance 0), which stores pages waiting to be processed. The next boxes represent pages discovered at increasing link distances (distance 1, then distance 2). This progression highlights the defining BFS principle: all pages at one depth are processed before the crawler advances to the next depth level. The visited-set box emphasizes that each URL is processed once, preventing duplicates and infinite loops.

\begin{figure}[htbp]
    \centering
    \begin{tabular}{@{}c@{\hspace{1.5em}}c@{\hspace{1.5em}}c@{\hspace{1.5em}}c@{}}
        \textbf{Seed URLs} & \textbf{Distance 0} & \textbf{Distance 1} & \textbf{Distance 2} \\
        \fbox{\parbox{3.4cm}{\centering Seed pages}} &
        \fbox{\parbox{3.4cm}{\centering Queue/Frontier}} &
        \fbox{\parbox{3.4cm}{\centering Pages in 1 hop}} &
        \fbox{\parbox{3.4cm}{\centering Pages in 2 hops}} \\
        $\rightarrow$ & $\rightarrow$ & $\rightarrow$ & \\
    \end{tabular}

    \vspace{0.8em}
    \fbox{\parbox{10.5cm}{\centering \small Visited set: already-processed URLs are not revisited, which prevents duplicates and keeps the crawl within the chosen depth limit.}}
    \caption{Breadth-First Search crawl strategy for a web spider. Starting from seed URLs, the crawler explores pages level by level in increasing distance from the seeds. The queue/frontier controls processing order, while the visited set prevents repeated processing of the same page.}
    \label{fig:bfs_spider}
\end{figure}

In practice, BFS means that the spider first visits all pages reachable in one link from the seeds, then all pages reachable in two links, and so on. This approach is well-suited for discovering distributed information (e.g., announcements, schedules, or event pages) because it expands uniformly across the site while still limiting exploration to a bounded depth. For thesis purposes, this figure clarifies why BFS is preferred for broad institutional discovery: it provides predictable coverage, transparent stopping criteria (maximum depth), and lower risk of missing near-root pages that often contain official updates.

\begin{table}[htbp]
\centering
\begin{tabular}{|p{105pt}|p{140pt}|p{160pt}|}
\hline
\textbf{Strategy} & \textbf{Core idea} & \textbf{Typical impact in institutional crawling} \\
\hline
BFS & Level-by-level expansion from seeds & Predictable coverage near root; good with max-depth constraints \\
\hline
DFS & Deep exploration along a path before backtracking & Faster deep discovery but may miss near-root categories under limited budget \\
\hline
Best-first (priority crawl) & Expand URLs by a score (keywords, patterns) & High relevance when scoring is accurate; risk of bias and missing diverse categories \\
\hline
\end{tabular}
\caption{Conceptual comparison of traversal policies used by web spiders.}
\label{tab:spider_traversal_strategies}
\end{table}

\subsection{Fetching, Politeness, and Crawl Constraints}
Autonomous extraction must operate responsibly and reliably. Even when crawling public pages, a spider should enforce rate limits, bounded concurrency, and timeouts to avoid stressing servers and to ensure predictable resource consumption. Constraints such as domain allowlists and maximum depth reduce accidental drift into unrelated sites. Deduplication (canonicalization of URLs, normalization of query strings, and content-based hashing) reduces wasted requests and limits data duplication. These constraints are not merely implementation details; they are foundational to turning a crawler into a maintainable system component \cite{Narayan2012}.

\subsection{HTML Parsing and Information Extraction}
After retrieval, HTML parsing transforms unstructured markup into analyzable representations. Parser pipelines generally include DOM construction, selector-based extraction, content cleaning, and field-level validation. Selector strategies may rely on structural cues (CSS classes, tag hierarchies), textual anchors (section titles, labels), or hybrid heuristics when pages are inconsistent. Because institutional websites evolve over time, parsers should be designed for resilience: they need fallback selectors, schema validation, and change-detection mechanisms that signal when extraction quality degrades \cite{Robbes2015}. Without such safeguards, autonomous extraction may silently produce incomplete or incorrect datasets.

\begin{table}[htbp]
\centering
\begin{tabular}{|p{110pt}|p{310pt}|}
\hline
\textbf{Pipeline stage} & \textbf{Purpose and typical outputs} \\
\hline
Discovery & Seed expansion, URL filtering, depth and domain constraints; outputs candidate URL set \\
\hline
Fetching & HTTP retrieval with timeouts, retries, caching; outputs raw HTML plus metadata (status, timestamp) \\
\hline
Parsing & DOM construction and normalization; outputs structured nodes suitable for selectors \\
\hline
Extraction & Field mapping (title/date/location/body); outputs typed records (e.g., \texttt{ScrapedEvent}) \\
\hline
Validation & Schema checks, date parsing, required fields, anomaly detection; outputs accepted/rejected records \\
\hline
Synchronization & Deduplication and upsert into platform storage; outputs versioned domain entities \\
\hline
\end{tabular}
\caption{A robust web extraction pipeline from discovery to synchronization.}
\label{tab:spider_pipeline}
\end{table}

\subsection{Quality, Change Detection, and Synchronization}
Extraction quality determines whether scraped content can be trusted as platform data. Common failure modes include missing fields, duplicate entries, incorrect date interpretations, and silent layout shifts in HTML. To mitigate these risks, spiders typically implement validation rules (required fields, range checks, format checks) and anomaly detection (unexpected drops in extracted item counts). Change detection is also critical: even when pages are stable, content changes should be detected efficiently to avoid unnecessary parsing and writes. Practical approaches include HTTP caching metadata (e.g., timestamps), content hashing, and record-level semantic identifiers so that updates are synchronized rather than duplicated \cite{Narayan2012}.

Finally, autonomy implies decision-making under uncertainty. A practical spider must determine which links deserve exploration, when to stop, and how to prioritize limited crawl budgets. BFS offers a transparent baseline policy, but it can be enhanced with relevance scoring functions and stop criteria that reflect platform needs. The core objective remains consistent: reduce manual data entry by turning distributed web content into reliable, structured, update-ready entities for platform services \cite{Sommerville2010}.